\section{证据使用与外推边界}
\label{sec:extension-evidence}

\subsection{四层证据不得混同}
解析证明建立恒等式、上限和下界；精确有理数核验固定反例；双精度数值检查验证实现与有限数据；独立系综实验验证生成模型、估计器与预测之间是否匹配。后一层通过不代替前一层证明，前一层成立也不保证任意真实影像流水线有匹配幅值或相同排序。

本轮保存的反例、参数族、逐条件误差界与测试均有可执行入口。实验说明列出随机种子、对象、灯和像素子集、噪声及先验约定。数值证书是解析证书的浮点评估，不称为采用定向舍入的计算机辅助证明。零更新、零任务和不可辨识方向按定义处理，未通过删去小的正边际、设置有利分母下限或丢弃失败条件来改善结果。

\subsection{历史证据与本轮新增的分离}
已有一般反例的 $14/109$、前轮实证 $R_A=0.55$ 及坐标匹配幅值 $53.744$ 都是从提供的证据包导入的历史事实；本轮没有将其冒充新测量。旧的 $R_A=0.90$、原候选下界和旧幅值记录仍可追溯，但不再担保当前推断。新结果单独登记，不回写先前实验文件。

源码中保留旧候选函数的数值行为，是复现兼容而非恢复其有效性。已有图像流水线的坐标、去规范和残差重抽样也未被静默改造成本轮匹配 GLS 实验；两者若要一一归因，必须继续构造同生成、同估计器、同投影及同基线的桥接。

\subsection{明确尚未完成的更强问题}

\begin{enumerate}
\item 物理反例允许逐灯正定耦合先验。共用同一对角先验、受限异方差权重或特定实际灯阵的更窄子类是否有额外强界，不能由该反例一并否定。
\item 固定两维、两候选、固定基线条件数的最坏值已给出；同时固定更新尺度、精化倍率、候选数与全部物理几何限制的联合最坏值，不在本轮已证范围内。
\item 两项近似的端点余项有普适二倍包围，但价值差可能发生谱项抵消。因此有限数据中界／实测误差小于十，不是该比例普遍有界的定理。判定精度时应使用绝对区间而非在真误差近零处分母不稳定的倍数。
\item 新 S1--S3 把受控系综的基线与有限前向偏差分开，并未把前轮 $53.744$ 完整分解成唯一的物理来源。新实验与历史估计器侧扰动存在结构差别，应继续保留这个负面解释边界。
\item 信息风险排序只对满足匹配、任务一致和可测裕量的设计对有保证。独立系统的有限验证不能升级为无条件跨系统泛化；真实角度误差也不自动等于加性反照率的二次任务。
\item 联合模型的上限和证书建立在同一固定线性化及商空间内。实物联合反照率--法向分配、非线性大扰动和改变阴影可见集后的全局保证，尚未由本轮合成验证完成。
\end{enumerate}

\subsection{可复用的研究结论}
本轮最直接可用于论文的方法补充，是把物理子类的反例、精确余项及其可算区间、匹配统计系综的条件风险关系和联合模型的固定空间结构命题各自写成独立结论。每个结论都可以在不改变历史实验数值的情况下加入正文或附录。保守或尚未闭合的部分应作为明确边界保留，而不是通过一个统一的正面口号掩盖。
